% =====================================================================
% Section I -- Introduction
% SpecAsync-UVM manuscript, draft 2 (revision of draft 1)
%
% CHANGES FROM DRAFT 1 -- all are number corrections or scope repairs,
% not rewrites of the argument:
%   (a) Pipelining ceiling 0.20-19.06%  ->  0.08-7.94%.  The old figure
%       came from a units mismatch (summed numerator over 5 trials,
%       single-trial denominator) that inflated every ceiling ~5x and
%       produced mathematically impossible >100% fractions on 4 of 6 T4
%       workloads.  See CEILING_BASIS_VERIFICATION.md / GATE_T3_REPORT.md.
%   (b) Hit rates 0.02% / 0.25%  ->  0.0132% / 0.2357% (marginal-diff of
%       pre-saturation reps; the old figures were computed off a saturated
%       ring).  Surrounding prose rewritten to state what the counter
%       actually measures.
%   (c) Prefetcher-disable cost 3.5x / 2.7-4.1x  ->  3.55x (T4), 3.88x
%       (5070 Ti) for stencil; +218% to +485% for cuFFT.
%   (d) "Six measurement artifacts"  ->  eleven, plus one candidate.
%   (e) \prov{} markers removed where the blocking work has since landed
%       (A1 worker telemetry, 5070 Ti cross-platform replication).  A3 is
%       addressed in prose rather than marked provisional, since B7's
%       12-core result largely answers it.
%
% NOT INCLUDED, deliberately: Gate B10 (C4 = prefetch ON + oracle) shows a
% large apparent benefit at iters=20 under oversubscription, but the effect
% oscillates in sign across iteration counts and a trace-ring phase-aliasing
% artifact has not yet been ruled out.  Nothing about B10 appears in this
% section until that is resolved.
% =====================================================================

\section{Introduction}

Unified Virtual Memory (UVM) lets GPU kernels address memory larger than device
capacity by migrating pages on demand, trading programmer effort for runtime
page-fault handling. That trade is not free: each GPU fault that misses in
device memory traps to the host driver, which must locate the faulting virtual
address range, decide a residency, construct and submit a migration, and update
page tables---all while holding a per-address-space lock. On fault-intensive
workloads this host-side path is a substantial fraction of execution time, and
it has accordingly attracted a decade of proposals to hide it: hardware and
software prefetchers, access-pattern predictors, tiling and scheduling
transformations, and speculative pre-staging of pages ahead of demand.

A large fraction of that literature is evaluated in simulation. Simulation is a
reasonable choice---modifying a production GPU driver is invasive, and
fault-path instrumentation at nanosecond resolution is difficult to obtain
otherwise---but it leaves a gap. A simulator models the costs its author chose
to model. If a proposed mechanism's benefit depends on a cost the simulator
abstracts away, or is defeated by an interaction with a component the simulator
does not implement, the proposal can appear sound and still fail on real
hardware. Closing that gap requires building the mechanism in a production
driver and measuring it.

This paper does that for one specific and widely-proposed idea: \emph{off-path
speculative page staging}. The mechanism is intuitive. When the driver services
a demand fault, it also predicts which page will fault next and hands that
prediction to an asynchronous worker outside the fault-service path. The worker
performs the residency lookup---and, optionally, the migration---concurrently
with the application's continued execution. If the prediction is correct, the
subsequent fault finds its work already done and the host-side cost is hidden.
Because the speculative work runs off the critical path, the mechanism appears
to be nearly free: a correct prediction saves time, and an incorrect one costs
only wasted background work.

We implemented this mechanism as \emph{SpecAsync-UVM}, a modification of
NVIDIA's open-source UVM kernel driver, and evaluated it on real hardware across
two GPU architectures, a range of workloads, and four prediction policies. Our
central finding is negative: \textbf{off-path speculative staging yields no
measurable wall-clock benefit, and the result is not attributable to poor
prediction.}

The evidence for that last clause is what distinguishes this study from a simple
null result. To separate prediction quality from achievable speedup, we
implement an \emph{oracle} policy that replays a previously recorded fault
trace, making prediction accuracy 100\% by construction. Under oracle prediction
with real page migration, and with every condition arranged in speculation's
favour, we observe no improvement over an unspeculated baseline on either of two
fault-intensive workloads: on GraphBFS the two are statistically
indistinguishable ($p=0.78$, Cohen's $d=0.10$, with a minimum detectable effect
of 0.16\% of baseline), and on a 2D stencil oracle speculation is measurably
\emph{slower} ($+6.7\%$, $p<10^{-6}$ Holm-corrected, $d=1.75$). The stencil
result reproduces on a second architecture in the same direction ($+3.60\%$
slower, $p<10^{-7}$). No predictor can stage pages more accurately than the
oracle; prediction accuracy is therefore not the bottleneck.

Three structural reasons emerge from our measurements. We report them as a
characterization rather than as a single mechanism, because the evidence
supports some more strongly than others.

\subsubsection*{Redundancy with the in-path prefetcher}
Off-path speculation is \emph{structurally redundant} with the driver's existing
in-path prefetcher. With the stock prefetcher enabled, speculative hit rates on
real workloads are near zero, because the prefetcher has already migrated the
pages that speculation would have predicted---it hides the very demand faults
speculation needs in order to act. Disabling the prefetcher exposes those
faults, but at a cost far exceeding anything speculation can recover:
$3.55\times$ on the stencil workload on a Tesla T4 and $3.88\times$ on an RTX
5070 Ti, and between $+218\%$ and $+485\%$ on cuFFT. That the two platforms
agree so closely despite an $8\times$ difference in interconnect bandwidth
indicates the prefetcher's value is not simply bandwidth-explained. Speculation
is squeezed between a regime where it has nothing to do and a regime whose entry
price it cannot repay.

\subsubsection*{Predictions are not consumed}
Even with the prefetcher disabled and prediction oracle-perfect, the driver
telemetry records speculative work being consumed at fault time only rarely:
0.0132\% and 0.2357\% of demand faults on our two T4 benchmarks, against
99.61\% for a deterministic synthetic probe under the same instrumentation. We
are deliberate about what this counter can and cannot establish. It increments
when a demand fault occurs and finds speculative work already staged for that
address, so it registers speculation that arrived in time to be \emph{useful}
but not in time to \emph{prevent} the fault; it bounds one mode of benefit
rather than measuring all of them. The negative result itself does not depend on
it---wall-clock time is measured directly and is independent of how hits are
counted. What the counter does establish is that the instrument works, and that
on real workloads predictions are overwhelmingly not consumed in the mode it can
see.

Our measurements narrow the cause to the latency of the asynchronous worker
relative to the arrival rate of demand faults. An uncontended round trip through
the worker costs 5.598\,\textmu s of dispatch latency against
0.300\,\textmu s of execution on the T4; under real workload load that
dispatch window inflates by three orders of magnitude ($1{,}685$--$1{,}894\times$),
while stencil faults arrive at roughly $188{,}000$ per second. The resulting race
margin is on the order of $10^3$. The blowup is load-dependent rather than a
fixed property of the design: it persists on a 12-core host, but at a
substantially lower magnitude ($432$--$438\times$), which indicates a queueing-pressure
threshold that additional cores raise without removing.

\subsubsection*{The headroom is small, and it is not where the mechanism aims}
The headroom available to \emph{any} such scheme is smaller than the mechanism's
framing suggests. Decomposing the fault-dispatch window into seven sub-phases,
we find that the overlappable portion---the work that could in principle be
moved off the critical path---bounds end-to-end speedup at between 0.08\% and
7.94\% across both platforms, driven almost entirely by how much of a workload's
wall-clock time is spent servicing faults at all. The dominant cost is servicing
under the address-space lock: 63--74\% of the dispatch window on the T4 and
52.1--76.9\% on the 5070 Ti. Lock \emph{wait}, by contrast, is negligible on both
($\approx 0.1\%$), so contention is not the problem. Source-level analysis of
the 595.71.05 tree shows the dominant phase is CPU-side residency bookkeeping
culminating in an asynchronous GPU submission rather than a blocking wait on the
GPU. It is therefore not latency that asynchrony can hide.

A cross-platform scaling comparison sharpens this. Moving the same workload from
a PCIe gen3 $\times$8 host to a gen5 $\times$16 host---an $8.00\times$ raw
interconnect improvement---yields $4.45\times$ on kernel-loop time,
$3.787\times$ on process wall-clock, and only $2.350\times$ on the fault-dispatch
window itself. The fault-servicing path benefits \emph{least} from a faster
interconnect of anything we measured, which suggests its cost floor is set by
per-fault CPU-side structure rather than by transport.

\subsubsection*{A premise of our own earlier work was wrong}
The mechanism was originally motivated by the belief that host-side metadata
discovery dominates fault-service cost. Corrected instrumentation shows metadata
accounts for 0.04\,\textmu s, under 1\% of that cost; dispatch and
migration dominate. The mechanism was aimed at the wrong target. We report this
directly because the correction is itself part of the paper's contribution: the
same instrumentation that refuted the premise is what makes the headroom bound
measurable.

\subsection{Contributions}

\begin{itemize}
  \item \textbf{An in-driver prototype on production hardware.} A working
  modification of NVIDIA's open-source UVM driver---under 400 net new lines,
  preserving the driver's safety invariants---implementing four prediction
  policies, two speculation depths, and debugfs telemetry, ported across driver
  versions and evaluated on a Tesla T4 (sm\_75, Turing, driver 595.71.05) and an
  RTX 5070 Ti (sm\_120, Blackwell, driver 595.84).

  \item \textbf{An oracle-based methodology separating prediction accuracy from
  achievable speedup.} By replaying recorded fault traces we establish an upper
  bound no heuristic predictor can exceed, allowing a negative result to be
  attributed to the mechanism rather than to prediction quality.

  \item \textbf{A rigorously validated negative result.} Measured under a
  pre-registered analysis plan with strictly interleaved run ordering,
  non-parametric testing, family-wise error correction, and reported minimum
  detectable effects---so that ``no benefit observed'' is a well-powered finding
  rather than an underpowered null.

  \item \textbf{A cross-architecture replication, with its evidential asymmetry
  stated.} Our two platforms differ along four axes simultaneously
  (architecture, driver version, host kernel, and virtualization), so a finding
  that reproduces across them is strong evidence of a driver-architectural
  property, while a finding that does not reproduce is causally uninterpretable.
  We label every result by which of these two categories it falls into rather
  than reporting all cross-platform comparisons as equivalent.

  \item \textbf{A quantified headroom bound and cost decomposition.} A
  seven-phase decomposition of the fault-dispatch window, with instrumentation
  overhead validated at 0.64\% (T4) and 0.331\% (5070 Ti, not statistically
  distinguishable from zero), yielding an end-to-end bound on what any off-path
  scheme could achieve, and a bandwidth-scaling ladder showing that the
  fault-service path is the component least improved by a faster interconnect.

  \item \textbf{A catalog of UVM measurement pitfalls.} Eleven distinct
  measurement artifacts encountered and corrected during this study, each
  presented with the false conclusion it would have produced---a reusable hazard
  list for anyone instrumenting this subsystem. Three of the eleven were caught
  not by the original analysis but by an independent verification pass over it,
  and each of those three inflated an absolute headline number while leaving a
  same-data ratio intact.
\end{itemize}

We do not claim that no form of speculation can benefit UVM. We claim,
specifically, that \emph{off-path asynchronous} speculation cannot, on this
driver architecture, for the reasons above---and that the headroom for
critical-path restructuring, while larger, is bounded and measurable. We close
by arguing that the productive direction indicated by our own decomposition is
in-path: reducing the cost of residency decision and migration construction
under the address-space lock, rather than predicting more accurately what to
fetch.
